\section{Data and Evidence Roles}
\label{sec:data}

\subsection{Cement activity and its period audit}

The official PBS Quantum Index of Manufacturing (QIM) reports
\SI{39.566}{\mega\tonne} of cement output for FY2023--24 \citep{pbsqim2024}. The State
Bank's FY2023--24 Statistical Supplement independently reproduces that value and provides a
five-year fiscal series \citep{sbpstatsupp2024}: 39.121, 49.797, 48.011, 41.448 and
39.566 Mt for FY2019--20 through FY2023--24. This establishes that the
\SI{48.011}{\mega\tonne} value previously attributed to FY2023--24 belongs to FY2021--22.
It is retained as a historical benchmark rather than a current structural alternative.

The annual core extends from FY2015--16 through FY2024--25 using held fiscal physical-production and trade records. The detailed reference-year analysis below remains FY2023--24. Publication revisions and source vintages are retained rather than interpreted as competing observations for one unchanged boundary.

QIM also reports FY2023--24 production of soda ash (\SI{0.785}{\mega\tonne}), billets and
ingots (\SI{4.914}{\mega\tonne}), glass plates and sheets
(\SI{20.205e6}{\metre\squared}) and lubricating oil
(\SI{130.513e6}{\litre}). These statistics establish the occurrence and scale of the
activities but do not provide every variable required by IPCC.
Glass needs mass, type, cullet and carbonate inputs; steel needs the production route and
carbon inputs; lubricants require non-energy consumption and oxidation rather than output.

Monthly UN Comtrade records for HS~252310 were summed from July 2023 through June 2024
\citep{uncomtrade2024clinker}. Exports total \SI{3.792074}{\mega\tonne} and imports
\SI{0.026131}{\mega\tonne}, with all twelve months present for both flows.

The panel now extends back to July 2014 for both flows, and every fiscal year of the
national series resolves to an observed trade term. Extending it required deciding what an
EMPTY month means, because the source returns no record both for a month in which nothing
was traded and for a month it has not processed, and the response does not distinguish
them. Pakistan's clinker imports are small and sporadic---\SI{0.026131}{\mega\tonne}
across the whole of FY2023--24 against \SI{3.792074}{\mega\tonne} exported---so empty
months are
common, and reading them as zeros would manufacture a complete panel from an incomplete
one, in the direction that lowers inferred domestic production. Each calendar year's
annual-frequency total was therefore retrieved as a control: a month with no record counts
as zero only where the monthly records of its calendar year already sum to that year's
annual total, so the source itself certifies that nothing further moved. Where they do not
sum, the month is unobtained and the fiscal years touching it are NE. Months lacking a record count as zero only after this compatible source-control check; unresolved coverage remains explicit. The
per-year verdicts are generated by \texttt{pkinv.sectors.ippu.trade\_evidence}. For the official
comparison year, January--June 2022 QIM observations were joined to July--December 2022 and
sum to \SI{44.331}{\mega\tonne}. Calendar-2022 WITS/UN Comtrade records give
\SI{1.341580}{\mega\tonne} of exports and \SI{0.003999}{\mega\tonne} of imports
\citep{witsclinker2022}.

\subsection{Producer and facility evidence}

We identified 16 cement producers relevant to FY2023--24. All sixteen public filings disclose
\SI{41.953170}{\mega\tonne} of clinker production
\citep{luckycement2024,dgkcement2024,attockcement2024,powercement2024,
cheratcement2024,pioneercement2024,thattacement2024,bestwaycement2024,
kohatcement2024,dandotcement2024,gharibwalcement2024,dewancement2025,
faujicement2024,mapleleafcement2024,flyingcement2024}.
Where disclosed, we also compiled opening and closing stocks, clinker purchases and imports,
clinker ground, domestic sales and exports. Missing flow terms remain unobserved, so company
balances are evaluated only where the necessary data are available.

Flying's scanned annual report gives exact FY2023--24 clinker and cement production of
\SI{0.322624}{\mega\tonne} and \SI{0.321500}{\mega\tonne}, respectively; Fecto's FY2025
annual report states the FY2023--24 comparative directly, \SI{0.651048}{\mega\tonne} of
clinker and \SI{0.724209}{\mega\tonne} of cement, where its earlier public briefing had
given only capacity and rounded utilisation
\citep{fectocement2024,fectocement2025,flyingcement2024}. All \num{16} census companies
therefore report an exact clinker quantity; what no company reports is the balance
around it (purchases, sales, grinding and stocks). The census is
\texttt{ippu\_cement\_producer\_census\_2024.csv}.

The facility registry identifies cement and fertiliser installations by operator, process route and administrative area. Coordinate provenance is retained for each record. The primary geographic package uses only locations compatible with its declared licence, while ODbL-derived locations retain a separate alternative. Company-wide clinker is split conditionally among works where plant-specific output is unavailable; unsupported placement remains unallocated. A point location does not itself establish plant production.

\subsection{Official reporting, chemicals and product-use data}

Pakistan's 2024 CRT submission contains a dedicated workbook for inventory year 2022
\citep{pakcrt2026}. We preserve ZIP and workbook hashes, worksheet names, cell addresses, raw
values and formulas. Table~2(I) provides category emissions; Table~2(I).A--H provides
activity, implied factors, native gases and recovery; and Table~2(II).B--Hs2 provides
species-level F-gas records.

Public company reports provide initial evidence for the next categories. Fatima
reports \SI{0.311242}{\mega\tonne} of calendar-2024 ammonia production at Sheikhupura; historical
records identify ammonia and nitric-acid output at Sadiqabad, dual-pressure nitric-acid
technology and CDM abatement projects at Sadiqabad and Pakarab \citep{fatima2024,pakarabcdm,
fatimacdm}. The NFDC annual review adds four extracted tables: domestic fertiliser production
by product and manufacturer, product-wise offtake split domestic and imported, plant-wise
fertiliser-sector gas, and the same gas monthly by plant for both fiscal years. The
feedstock/fuel split is therefore observed for the reference year itself---
\SI{83.5}{\percent}/\SI{16.5}{\percent} in 2023--24, and
\SI{83.2}{\percent}/\SI{16.8}{\percent} in 2024--25---rather than borrowed backwards from
the later year. The reference year's total, \SI{751.9}{} MMCFD, agrees with the
regulator's independently published \SI{756}{} to \SI{0.5}{\percent}. These
support a four-pathway activity reconciliation for ammonia---official filing, plant census,
product stoichiometry and gas balance---whose pathways do not yet agree. National plant
coverage, recovered \COtwo{}, and matched abatement utilisation
remain incomplete, so these observations support the category assessment but are not used to
estimate national emissions. Calendar-2022
and calendar-2024 records are audited separately. Verified CDM reductions are retained by
monitoring period and are not carried forward as current \NtwoO{} emissions or abatement rates.
For Pakarab, the period-specific audit preserves 2008 line production, operating days,
nameplate capacity, tertiary De-\NtwoO{} technology and the approved baseline cap without
turning those historical observations into FY2023--24 activity.

Pakistan's Ozone Secretariat profile reports HFC consumption of
\SI{9.815266}{\mega\tonne} \COe{} in 2024 \citep{pakistanozoneprofile}; Pakistan ratified
the Kigali Amendment on 22~October~2025 as an Article~5 Group~2 party, whose baseline is the
2024--2026 consumption average, and the profile records that no HFC licensing system is yet
established. Consumption is not
annual leakage. Estimation therefore requires annual data by species and application on
imports, exports, production, destruction, new and servicing charge, equipment retirement and
recovery, together with the parameters needed by the equipment-bank model. Two
free instruments sit beside the consumption anchor: UN~Comtrade's HS-2022 HFC lines
(net imports of \SI{6.0}{\kilo\tonne} of substance in 2024, bulk plus mixtures --- a
customs corroboration of supply that cannot resolve species) and a one-stratum
\COe{}-weighted bank model run on the consumption series itself under declared IPCC
default losses, whose 2024 reading of \SI{4.42}{\mega\tonne} (2.3--6.2 across the
5--15\,\% loss-rate cases, not a bound) ships as the screening configuration
\texttt{scope\_plus\_fgas\_bounded\_bank}, outside every reference total.

\begin{table*}[htbp]
\centering
\footnotesize
\begin{tabular}{@{}p{0.13\linewidth}p{0.20\linewidth}p{0.25\linewidth}p{0.27\linewidth}@{}}
\toprule
\textbf{Category} & \textbf{FY2023--24 status} & \textbf{Available public evidence} & \textbf{Evidence still required} \\
\midrule
2.A.1 Cement & E & QIM, trade, 16 producer filings & common boundary, plant splits, clinker movements, stocks and chemistry \\
2.A.2 Lime & NE & Mineral/industry screen & Marketed and captive lime by type \\
2.A.3 Glass & NE & QIM area output & Mass, glass type, cullet and carbonate inputs \\
2.B.1 Ammonia & E & NFDC production by manufacturer, plant gas feed/fuel split, one FY2024 plant, CRT/history & Plant-level output and recovered \COtwo{} for the Tier~2 step \\
2.B.2 Nitric acid & NE & Plant and CDM evidence & FY2024 100\%-acid output and abatement operation \\
2.C.1 Iron/steel & E & QIM billets/ingots & Process route and carbon inputs \\
2.D.1 Lubricants & E & QIM output & Non-energy use and oxidation \\
2.F HFCs & NE & Ozone consumption & Species/application flows and equipment banks \\
2.G SF$_6$/other & NE & Occurrence screen & Equipment banks, refills, recovery and use data \\
\bottomrule
\end{tabular}
\caption{Evidence status for the principal IPPU categories. The complete supplementary
matrix covers every detailed 2.A--2.G category and documents the data needed for estimation.}
\label{tab:completeness}
\end{table*}
